\section{Post-Pause Experiment Blueprint}

This section turns the evaluation plan into a concrete, execution-ready blueprint. It is still a plan, not a reported result set, because the experiment pause remains active.

\subsection{Measurement Goals}

The first goal is to establish a reliable recoverability baseline. Under synchronized settings, we need to know not only overall decode success rate but also how success changes with message length, prompt class, and language.

The second goal is to characterize runtime cost. Encode and decode latency, token throughput, and token consumption are needed for practical deployment decisions.

The third goal is to measure embedding utilization in a way that can be compared with prior work, while remaining faithful to Ghostext's packetized protocol.

The fourth goal is to collect security-facing indicators under the passive-censor model, including distributional deviation and preliminary detector outcomes.

\subsection{Prompt and Message Regimes}

We will evaluate multiple prompt regimes instead of a single template family. At minimum, prompts should cover narrative prose, informational style, and conversational style. Since Ghostext targets bilingual usability, both English and Chinese prompts should be included.

Message regimes should include short, medium, and long payload bands. Because Ghostext uses authenticated packetization, payload size impacts both packet bytes and required token budget; reporting should therefore separate plaintext size, packet size, and consumed token length.

\subsection{Configuration Grid}

The grid should vary only a small number of parameters at once to keep attribution clear. A practical initial grid is:

\begin{itemize}
\item candidate truncation parameters (top-$p$, max-$k$),
\item quantization total frequency,
\item segment token budgets,
\item retry and entropy guard parameters.
\end{itemize}

All other settings should remain fixed within each sweep, including model file identity and backend build.

\subsection{Recoverability Reporting}

Recoverability should be reported with two complementary views.

The first view is aggregate success percentage over trials.

The second view is failure decomposition by explicit error class. Because Ghostext exposes class-specific failures (integrity mismatch, synchronization mismatch, low-entropy retry exhaustion, token-budget exhaustion, and others), decomposition is both feasible and informative.

This decomposition is critical for engineering actionability: different failure classes imply different mitigations.

\subsection{Runtime and Cost Reporting}

Runtime should include at least wall-clock encode latency, decode latency, and tokens-per-second metrics. Cost should include packet-carrying token count and total output token count, because optional natural tails affect total transmission length without increasing payload.

For fair comparison across settings, report median and dispersion statistics rather than single-point means only.

\subsection{Utilization Definitions}

Ghostext should report two utilization metrics:

\textbf{Packet efficiency}: embedded packet bits per packet-resolving token.

\textbf{Transmission efficiency}: embedded packet bits per total transmitted token.

Both are useful. Packet efficiency reflects core embedding mechanics, while transmission efficiency captures user-visible communication cost.

\subsection{Distribution-Facing Security Indicators}

Following NLS-style framing~\cite{ziegler2019neural}, we will compare stego outputs with baseline model outputs under matched prompt families and report divergence-oriented indicators. Because finite truncation and quantization are part of Ghostext, the interpretation should explicitly mention parameter settings for each reported number.

The goal of this block is not to claim perfect indistinguishability but to map where divergence remains small enough to support the passive-censor narrative.

\subsection{Detector-Facing Preliminary Block}

Detector experiments will be labeled preliminary until robustly replicated. The planned block includes representative classical/neural steganalysis baselines and one LLM-based evaluator as a human-like proxy, consistent with recent practice discussed in OD-Stega~\cite{huang2026odstega}.

For this block, reproducibility metadata is especially important because detector performance is sensitive to training data composition and prompt distribution shift.

\subsection{Artifact Packaging}

Each result bundle should include:

\begin{itemize}
\item raw JSON summaries,
\item scripts and command lines,
\item exact implementation commit hash,
\item paper commit hash,
\item environment manifest (hardware, OS, dependency versions).
\end{itemize}

This package should be sufficient for a teammate to regenerate tables and plots from scratch.

\subsection{Decision Criteria for Next Iteration}

After first post-pause runs, we should decide whether to prioritize optimization integration or robustness extensions. A practical decision rule is:

If recoverability and reproducibility are strong but utilization is low, prioritize OD-style distribution optimization integration.

If recoverability failures cluster in mismatch classes, prioritize synchronization hardening and tooling.

If detector results degrade quickly with payload pressure, prioritize adaptive policy tuning and stronger security-focused evaluation before adding new features.

\subsection{Expected Outcome of This Blueprint}

The intended outcome is a measured, traceable transition from design-intent claims to evidence-backed claims. Once this blueprint is executed, the manuscript can replace conservative placeholders with concrete quantitative statements while keeping the same claim-discipline structure.
